\chapter{Uvod}
\label{ch:uvod}

% ------------------------------------------------------------
\section{Problem}
\label{sec:problem}

Od 2026.\ obveznik uz dostavu eRačuna primatelju mora prijaviti i podatke o
njemu Poreznoj upravi. Budući da su to odvojeni postupci, jedan može uspjeti,
a drugi ne. Ako na fiskalizacijski zahtjev ne stigne
odgovor, obveznik ne zna je li zahtjev obrađen. U razmatranom primjeru
obveznik ponovno šalje zahtjev, a Sustav za fiskalizaciju odgovara šifrom
S008. Ta šifra označava da zapis s istim identifikatorom već postoji.

Tumačenje šifre S008 otežano je kod ispravka podatka koji ne utječe na obračun
poreza. Takav se ispravak izdaje pod istim brojem računa i s indikatorom
kopije, ali indikator nije dio složenog identifikatora koji se koristi u
provjeri S008. Ako ispravak zadrži datum izvornika, identifikator
ostaje isti i opisana provjera vodi u odbijanje. Ako dobije novi datum i
referencu na prethodni račun, identifikator se mijenja i ta se provjera
izbjegava. Izvori ne određuju koji je datum ispravka mjerodavan.
Prihvatljivost varijante s novim datumom u stvarnom sustavu nije provjerena.
Nalaz se zato odnosi na nedovoljno određenu semantiku sučelja.

Rad iz propisanog tijeka izvodi jedanaest načina otkazivanja, među kojima je
i opisani slučaj. Dio ograničenja proizlazi iz svojstava raspodijeljene
razmjene, poput nemogućnosti da pošiljatelj iz samog izostanka odgovora
utvrdi ishod zahtjeva. Drugi dio proizlazi iz obvezne arhitekture i pravila
koja ne određuju semantiku kvara. Za tu skupinu rad upotrebljava radni pojam
\emph{regulatorno inducirano ograničenje}.

Literatura opisuje mehanizme za ublažavanje ograničenja raspodijeljene
razmjene. Kada ograničenje proizlazi iz propisa, obrasci mogu suziti prozor
nepoznatog ishoda, izričito zabilježiti takvo stanje i sačuvati dokaz o
postupanju. Ne mogu, međutim, odrediti pravno značenje pojma koje u propisu
nije određeno. Poglavlje~\ref{ch:rasprava} zato učinak obrazaca ocjenjuje
uvjetno, ovisno o granici sustava na kojoj se primjenjuju i njihovoj
konfiguraciji.

% ------------------------------------------------------------
\section{Ciljevi i istraživačka pitanja}
\label{sec:ciljevi}

Cilj rada je utvrditi koji se načini otkazivanja mogu izvesti iz propisanog
tijeka fiskalizacije eRačuna i koje od njih pokrivaju obrasci tolerancije na
greške. Pritom se razmatra kako primjena obrazaca utječe na složenost izvedbe,
kašnjenje i rizik prekoračenja propisanog roka. Analiza se temelji na
objavljenim dokumentima, pa se izvedene tvrdnje mogu provjeriti bez pristupa
stvarnom sustavu.

Rad odgovara na tri pitanja.

\begin{description}[leftmargin=1.6cm, style=nextline]
  \item[RQ1] Koji se načini otkazivanja mogu izvesti iz propisanog tijeka
  fiskalizacije eRačuna i u kojim završnim stanjima ostavljaju sustav s obzirom na
  zakonsku obvezu?

  \item[RQ2] Koji obrasci tolerancije na greške pokrivaju koja od izvedenih
  završnih stanja, koja ostaju samo detektirana ili nepokrivena te proizlazi li
  granica pokrivenosti iz svojstava raspodijeljenih sustava ili iz propisa?

  \item[RQ3] Koju cijenu u složenosti, kašnjenju i riziku prekoračenja
  propisanog roka nose ti obrasci i u kojim slučajevima pogoršavaju usklađenost
  umjesto da je poboljšaju?
\end{description}

Drugim se pitanjem ispituje i proizlaze li ograničenja obrazaca iz propisa.
Time se provjerava teza rada, uz mogućnost da analiza pokaže potpunu
pokrivenost ili ograničenje koje proizlazi iz propisa.

Rezultati rada obuhvaćaju katalog načina otkazivanja, radni pojam regulatorno
induciranog ograničenja i matricu pokrivenosti. Katalog se temelji na
propisanim obvezama, pa njegova primjena nije vezana uz jedan sustav. Ipak,
ne pojavljuju se svi slučajevi u svakoj arhitekturi. Obveznik koji je vlastita
pristupna točka nema granicu prema posredniku opisanu slučajevima D2 i E1,
dok su za obveznika koji koristi posrednika relevantna oba slučaja. Skupine
A i C proizlaze iz obveza koje ne ovise o tom izboru.

Pojam regulatorno induciranog ograničenja povezuje slučajeve u kojima
ograničenje proizlazi iz propisane arhitekture ili nedovoljno određenih
pravila. Matrica pokrivenosti daje podlogu za izbor obrazaca, uz navedene
uvjete njihove primjene. Obuhvaća i slučajeve u kojima učinak ovisi o
konfiguraciji i preostalom roku.

Istek vremena, ponavljanje, prekidač i saga razmatraju se kao mogući obrasci
za primjenu u predloženom rješenju. Za svaki se obrazac ispituju uvjeti
primjenjivosti, a odluka o njegovu uključivanju ili izostavljanju obrazlaže
se rezultatima analize.

% ------------------------------------------------------------
\section{Opseg i struktura}
\label{sec:opseg-rada}

Studija slučaja obuhvaća arhitekturu, stanja i odgovornosti sudionika u
cijelom propisanom tijeku. Predmet analize je referentna arhitektura izvedena
iz propisa, a ne postojeći sustav. Izvršni prototip modelira dvije granice
na kojima se pojavljuju slučajevi B1, B2, D1 i E1. Na njima uspoređuje dvije
izvedbe: osnovna stvara zapis predaje tek nakon primitka odgovora, bez
prethodnog zapisa namjere, a predložena zapisuje namjeru prije slanja.

Evaluacija povezuje analitičku matricu s četiri izvršna scenarija.
Učinkovitost se procjenjuje prema tome koliko je stanje sustava nakon istog
kvara poznato i dokazivo. Propusnost i vrijeme odziva nisu mjerila ove
evaluacije. Njezin opseg i ograničenja obrazlaže
poglavlje~\ref{ch:rasprava}.

Poglavlje~\ref{ch:pozadina} opisuje pravni i tehnički kontekst.
Poglavlja~\ref{ch:raspodijeljeni} do~\ref{ch:tijekovi} objašnjavaju djelomični
otkaz i klasifikaciju otkaza te obrasce tolerancije na greške i granice
transakcije u raspodijeljenim tijekovima. Poglavlje~\ref{ch:studija} izvodi
načine otkazivanja i odgovara na RQ1. Poglavlje~\ref{ch:implementacija}
opisuje ilustrativni prototip, a poglavlje~\ref{ch:evaluacija} matricom
pokrivenosti i izvršnim scenarijima odgovara na RQ2.
Poglavlje~\ref{ch:rasprava} razmatra troškove i ograničenja primjene obrazaca
te odgovara na RQ3.
